\documentclass{article}
\usepackage[utf8]{inputenc}
\usepackage[T1]{fontenc}
\usepackage[a4paper]{geometry}		% Reduire les marges
\geometry{hmargin=3cm,vmargin=3.5cm}
\usepackage[english]{babel} 
\usepackage{amsfonts}
\usepackage{amsthm}
\usepackage{amsmath}
\usepackage{amssymb}
\usepackage{hyperref}
\usepackage{array}
\usepackage[svgnames]{xcolor}
\usepackage{xpatch}
\usepackage{tikz}
\usepackage{mathrsfs}
%\usepackage{graphicx}
%\usepackage{caption}
\usepackage[svgnames]{xcolor}
\usepackage{eurosym}
\usetikzlibrary{positioning,calc}
\renewcommand{\thesection}{\Roman{section}}
\xpatchcmd{\proof}{\itshape}{\normalfont\proofnameformat}{}{}
\newcommand{\proofnameformat}{\bfseries}
\theoremstyle{definition}
\newtheorem{defi}{Definition}[section]
\newtheorem{theo}{Theorem}[section]
\newtheorem{prop}[theo]{Proposition}
\newtheorem{lem}[theo]{Lemma}
\newtheorem{coro}[theo]{Corollary}
\newtheorem{exe}{Example}[section]
\newtheorem{rem}{Remark}[section]
\newtheorem{nota}{Notation}[section]
\renewcommand\qedsymbol{$\blacksquare$}
\numberwithin{equation}{section}

\DeclareMathOperator{\R}{\mathbb{R}}
\DeclareMathOperator{\Z}{\mathbb{Z}}
\DeclareMathOperator{\Q}{\mathbb{Q}}
\DeclareMathOperator{\N}{\mathbb{N}}
\DeclareMathOperator{\E}{\mathbb{E}}
\DeclareMathOperator{\C}{\mathcal{C}}

\title{Exo 8, corrigé}       
\author{Guillaume Woessner}
\date{}                       % sans date : celle du jour de compilation

%\pagestyle{headings}          % Pour mettre des entetes avec les titres
                              % des sections en haut de page
\sloppy                       % Ne pas faire deborder les lignes dans la marge

\begin{document}

\maketitle                    % Faire un titre utilisant les donnees
                              % passees : \title, \author et \date

\begin{abstract}
\begin{center}
Les familles U.I. et les processus de branchement.
\end{center}
\end{abstract}

%\tableofcontents              % Table des matieres


Dans tout ce document, l'espace de probabilité filtré considéré est $(\Omega,\mathcal{F},(\mathcal{F}_n)_n,\mathbb{P})$, et on supposera que $(X_n)_n$ est une martingale adaptée à la filtration.




%\paragraph{Exercice 1}
%Montrez que la famille $(X_n)_n$ est UI si et seulement si les deux conditions suivantes sont vérifiées,
%\begin{itemize}
 %   \item[$\bullet$] $\sup_n \E|X_n| < \infty $
  %  \item[$\bullet$] $\forall \varepsilon >0,~\exists \delta>0,~\forall B\in\mathcal{F},~\mathbb{P}(B)<\delta \rightarrow \sup_n\E(|X_n\mathbf{1}_{B})<\varepsilon. $
%\end{itemize}


\paragraph{Exercice 1}
Supposons qu'il existe un $p>1$ tel que $\sup_n\E|X_n|^p<\infty$.\\
Montrez que $(X_n)_n$ est UI.

\begin{proof}
 On peut calculer que 
 \begin{align*}
 \E[|X_n|\mathbf{1}_{|X_n|>K_\varepsilon}] &= \E\left[|X_n|^p \dfrac{\mathbf{1}_{|X_n|>K_\varepsilon}}{|X_n|^{1-p}}\right]\leq \E\left[|X_n|^p \dfrac{\mathbf{1}_{|X_n|>K_\varepsilon}}{K_\varepsilon^{1-p}}\right]\\
 &=  \dfrac{1}{K_\varepsilon^{1-p}}\E[|X_n|^p\mathbf{1}_{|X_n|>K_\varepsilon}] \leq \dfrac{\sup_n\E|X_n|^p}{K_\varepsilon^{1-p}} 
 \end{align*}
 Ainsi choisir $K_\varepsilon$ de telle sorte que $\varepsilon= \dfrac{\sup_n\E|X_n|^p}{K_\varepsilon^{1-p}} $ fait l'affaire.
\end{proof}


%\paragraph{Exercice 3} Supposons que $\sup_n |X_n|<Y$ où $Y$ est intégrable.\\
%Montrez que $(X_n)_n$ est UI.


\paragraph{Exercice 2} Soit $X$ une variable aléatoire intégrable.\\
Montrez que la famille $(\E[X~|~\mathcal{G}]~:~\mathcal{G}\subset \mathcal{F}$ est une sous-tribu), est UI.\\
\textit{Indication} : Vous pouvez utiliser le lemme suivant :\\
{\small Si $X$ est intégrable, alors il existe $h~:~\R^+\to \R^+$ une fonction continue, convexe et croissante vérifiant $\lim_{x\mapsto \infty} \frac{h(x)}{x}=\infty$ telle que $\E[h(X)]<\infty$.}

\begin{proof}
 Soit $h$ donnée par le Lemme. Ainsi on calcule
 $$ \E[h(\E[X~|~\mathcal{G}])] \leq \E[\E[h(X)~|~\mathcal{G}]] = \E[h(X)] < \infty. $$
 Un raisonnement analogue à celui de l'Exercice 1 permet alors de conclure.
\end{proof}
\begin{proof}[Preuve du Lemme]
 Pour tout $n\in\N$ on sait qu'il existe $a_n\in\R^+$ tel que $\E[|X|\mathbf{1}_{|X|>a_n}]\leq \frac{1}{n^2}$, et on peut supposer sans perte de généralité que $(a_n)$ forme une suite croissante qui tend vers l'infini.\\
 On pose alors 
 $$g(x) := \sum_n |x| \mathbf{1}_{|x|>a_n} \qquad \text{et} \qquad h(x):=\sum_n (x-a_n)^+.$$
 On remarque facilement que pour tout $x\in\R^+$, on a $h(x)\leq g(x) <\infty$. De plus, étant une somme de fonctions convexes continues et croissantes, $h$ est également convexe continue et croissante. Enfin on peut aussi montrer que $\lim_{x\mapsto \infty} \frac{h(x)}{x}=\infty$.\\
 Alors on calcule
$$
\E[h(|X|)] \leq \E[g(|X|)] = \E\left[ \sum_n |X|\mathbf{1}_{|X|>a_n} \right] = \sum_n \E[|X|\mathbf{1}_{|X|>a_n}]\leq \sum_n \dfrac{1}{n^2}<\infty.
$$
\end{proof}


\paragraph{Exercice 3} Preuve alternative du fait que la population s'éteint dans le cas $R < 1$.
\begin{itemize}
    \item Revisitez la preuve de convergence p.s. des martingales, et vérifiez que $\E[X_\infty] \leq \liminf \E[X_n]$
    \item Dans la preuve du théorème de Galton-Watson, montrez que (avec les notations de la preuve), si $R<1$, alors on a $\lim_n\E[X_n]$, $\E[X_\infty]=0$ et $X_\infty =0$. 
\end{itemize}

\begin{proof}
 Pour le premier point il suffit d'utiliser le lemme de Fatou.\\
 Pour le second point on remarque que $M_n=\frac{X_n}{R^n}$ est une martingale et que par conséquent 
 $$ \E[X_n] = R^n\E[M_n] = R^n \E[M_0] = R^n \mapsto 0. $$
 On conclut en utilisant le premier point.
\end{proof}


\paragraph{Exercice 4} Fin de la preuve du théorème de Galton-Watson dans le cas $R > 1$.\\
On veut montrer que la probabilité d'extinction (notée $1-p$ dans le cours) est égale à la plus petite solution $q_0$ de $g(z) = z$, où $g(z) = \E(z^Y)$ où $Y$ est la loi de régénération pour un processus de branchement. Pour ce faire
\begin{itemize}
    \item Montrez que $\mathbb{P}(X_n = 0) = g(g \dots g(0))$ où on applique $g$ exactement $n$ fois, e.g. $\mathbb{P}(X_1 = 0) = g(0)$ et $\mathbb{P}(X_2 = 1) = g(g(0))$ etc...
    \item Argumentez que $1-p = \lim_{n \to \infty} \mathbb{P}(X_n = 0)$
    \item En utilisant la monotonicité de $g$, montrez que $1-p = q_0 < 1$.
\end{itemize}

\begin{proof} On procède comme suit :
\begin{itemize}
    \item On raisonne par récurrence. Les cas $n=0$ et $n=1$ sont évidents, et on suppose alors qu'il existe $n\in\N$ tel que $\mathbb{P}(X_n=0)=g^n(0)$. On calcule alors
    \begin{align*} 
    \mathbb{P}(X_{n+1}=0) &= \sum_k \mathbb{P}(X_{n+1}=0~|~X_1=k)\mathbb{P}(X_1=k)\\
    &=^\star \sum_k \mathbb{P}(X_n=0)^k\mathbb{P}(Y=k)\\
    &= \sum_k g^n(0)^k \mathbb{P}(Y=k)\\
    &= g\circ g^n(0) = g^{n+1}(0). 
    \end{align*}
    \item Comme $(\{X_n=0\})_n$ est une famille croissante d'évènements, on peut écrire que 
    $$ \lim_n\mathbb{P}(X_n=0) = \mathbb{P}(\cup_n\{X_n=0\}) = 1-p. $$
    \item On a obtenu que $1-p=\lim_n\mathbb{P}(X_n)=0 = \lim_n g^n(0)$.\\
    Ceci définit une suite déterminée par une relation de récurrence. Comme on sait que $g(0)=\mathbb{P}(Y=0)\geq 0$ et que $g(1)=\E[1^Y]=1$, en utilisant les propriétés des suites déterminées par des relations de récurrence on sait que $1-p$ est égal à la plus petite solution $q_0$ de $g(z)=z$. 
\end{itemize}
$\star$ cette égalité est vraie car on peut montrer facilement que conditionnellement à $\{X_n=k\}$, chacune des $k$ "lignées" est indépendante des autres et a la même loi qu'un processus de branchement.
\end{proof}



\paragraph{Exercice 5} Correction Exercise 7 feuille 6.
Vérifiez bien que la convergence a lieu \textit{p.s.} et dans $L^1$ (en fait dans $L^p$ avec $p < 2$).
D'autre coté, trouver une sous-martingale telle que la convergence n'ait pas lieu dans $L^2$ (indication: similaire à la solution de l'exo 3 feuille 7).

\begin{proof}
On peut poser $X_n~:~[0,1]\to\R$ tel que $X_n(\omega)=-2^{n/2}\mathbf{1}_{[0,2^{-n}[}(\omega)$.
\end{proof}


\paragraph{Exercice 6}
Montrez que toute sous-martingale positive d'espérance constante est une martingale.
\begin{proof}
On sait que $\E[M_{n+1}-M_n~|~\mathcal{F}_n]\geq 0$. D'autre part on sait que $$ \E[\E[M_{n+1}-M_n~|~\mathcal{F}_n]]=\E[M_{n+1}-M_n] = \E[M_{n+1}]-\E[M_n] = 0. $$
Ainsi on a que $\E[M_{n+1}-M_n~|~\mathcal{F}_n]\geq 0$ presque sûrement (voir Exerice 1 Feuille 7).
\end{proof}

\end{document}













